\subsection{自同态的行列式}

设 M 是一个具有 n 个元素的有限基的 A-模，u 是 M 的一个自同态。A-模 \(\bigwedge^{n}(\mathbf{M})\) 是一个单生成自由模，即与A同构（第8号，定理1的推论1）。 \(\bigwedge^{n}(u)\) 是该模的一个自同态，因此是一个位似 \(z \mapsto \lambda z\) ，其位似比 \(\lambda \in A\) 被唯一确定（II，§2，第 3 号，命题 5）。

定义 1. 一个自同态 \(u\) （自由 A-模 \(M\) 的、维数 \(n\) 为有限者；II，§7，第 2 号，命题 3 的推论及注记 1）的行列式，记作 \(\det u\) ，是标量 \(\lambda\) （它使得 \(\bigwedge^{n}(u)\) 是比率为 \(\lambda\) 的位似）.

根据 §7 第 2 号的公式 (4)，det u 是使得下式成立的唯一标量：

\[
u \left(x _ {1}\right) \wedge u \left(x _ {2}\right) \wedge \dots \wedge u \left(x _ {n}\right) = (\det u). x _ {1} \wedge x _ {2} \wedge \dots \wedge x _ {n}\tag{1}
\]

对每个序列 \((x_{i})_{1\leqslant i\leqslant n}\) （ \(n\) 个 M 的元素）成立。如果 \(\det (u) = 1\) ，则称 \(u\) 为幺模的。

命题1.（i）若 \(u\) 和 \(v\) 是两个有限维自由A-模M的自同态，则

\[
\det (u \circ v) = (\det u) (\det v).\tag{2}
\]

\begin{enumerate}
\def\labelenumi{(\roman{enumi})}
\setcounter{enumi}{1}
\tightlist
\item
  \(\operatorname{det}(1_{\mathbf{M}}) = 1\) ；对每个自同构 \(u\) （ \(\mathbf{M}\) 的自同构）， \(\operatorname{det} u\) 在 \(\mathbf{A}\) 中可逆，且
\end{enumerate}

\[
\det (u ^ {- 1}) = (\det u) ^ {- 1}.\tag{3}
\]

若 \(n\) 是 \(\mathbf{M}\) 的维数，则这由关系式 \(\bigwedge^{n}(u \circ v) = (\bigwedge^{n}(u)) \circ (\bigwedge^{n}(v))\) （§ 7，第 2 号，公式 (3)）直接得出。

设 M 为具有有限基 \((e_{i})_{1\leqslant i\leqslant n}\) 的自由 A-模；给定 M 的 n 个元素组成的一个序列 \((x_{i})_{1\leqslant i\leqslant n}\) ，此序列关于给定基 \((e_i)\) 的行列式（当基不致引起混淆时记作 \(\det(x_1, x_2, \ldots, x_n)\) ）是 M 的自同态 \(u\) （它满足 \(u(e_i) = x_i\) ， \(1 \leqslant i \leqslant n\) ）的行列式。于是由公式 (1)

\[
x _ {1} \wedge x _ {2} \wedge \dots \wedge x _ {n} = \det (x _ {1}, x _ {2}, \dots , x _ {n}) e _ {1} \wedge e _ {2} \wedge \dots \wedge e _ {n}\tag{4}
\]

且此关系刻画了映射 \((x_{i})\mapsto\det(x_{1},x_{2},\ldots,x_{n})\) （从 \(\mathbf{M}^{n}\) 到 A）。它表明该映射是一个交错 \(n\) -线性形式。由于根据 §7 第 4 号命题 7，交错 \(n\) -线性形式组成的 A-模典范同构于 \(\bigwedge^{n}(\mathbf{M})\) 的对偶，且 \(\bigwedge^{n}(\mathbf{M})\) 同构于 A，由此可见每个交错 \(n\) -线性形式在 \(\mathbf{M}^{n}\) 上都可以写成

\[
(x _ {1}, \dots , x _ {n}) \mapsto \alpha \det (x _ {1}, x _ {2}, \dots , x _ {n})
\]

对于某个 \(\alpha \in \mathbf{A}\) .

命题2. 设 M 是一个具有有限基 \((e_i)_{1 \leqslant i \leqslant n}\) 的自由 A-模， \(v\) 是 M 的一个自同态。对每个序列 \((x_i)_{1 \leqslant i \leqslant n}\) （由 \(n\) 个 M 的元素组成），

\[
\det (v (x _ {1}), \dots , v (x _ {n})) = (\det v) \det (x _ {1}, \dots , x _ {n}).\tag{5}
\]

若 \(u\) 是 \(M\) 的、使得 \(u(e_i) = x_i\) 对每个 \(i\) 成立的自同态，则

\[
v (x _ {i}) = (v \circ u) (e _ {i})
\]

因此，(5) 由 (2) 得出。

